\documentclass[11pt,a4paper]{article}
\usepackage[margin=2.6cm]{geometry}
\usepackage{amsmath,amssymb,amsthm,mathtools,microtype,booktabs,array,tabularx}
\usepackage[shortlabels]{enumitem}
\usepackage[colorlinks=true,linkcolor=black,urlcolor=blue]{hyperref}
\hypersetup{pdftitle={A human proof of the one-linear-gate case},pdfauthor={AI-written draft, pending human review}}
\newtheorem{theorem}{Theorem}
\newtheorem{lemma}[theorem]{Lemma}
\newtheorem{proposition}[theorem]{Proposition}
\newcommand{\F}{\mathbb F_2}
\newcommand{\Aff}{\mathrm{Aff}}
\newcommand{\Span}{\operatorname{span}}
\newcommand{\rk}{\operatorname{rank}}
\newcommand{\supp}{\operatorname{supp}}
\newcommand{\status}[1]{\hfill\textsf{[#1]}}
\setlength{\parskip}{3pt}
\title{A human proof of the one-linear-gate case\\of the low-half multiplier}
\author{AI-written draft, pending human review}
\date{1 November 2026}
\begin{document}
\maketitle
\begin{abstract}
We give a human proof that no five-AND circuit with one affine-input gate computes the three low bits of the three-bit product. The quadratic first nonlinear gate reduces to the proved two-linear-gate theorem. For a cubic first gate, a core-supported case follows from a radical-pairing extension of the coset obstruction. In the noncore case, parity on two core hyperplanes excludes ordinary products, while a multiplier lemma forces every feedback multiplier to equal one modulo $(a_0b_0)$. The ideal-extraction lemma then supplies the final rank contradiction. Every new lemma is backed by an exhaustive check of its application domain, including 171,966,464 multiplier pairs. An independent full circuit search also passes. The draft awaits human review.
\end{abstract}

\section{Setting and imports}
The variables are $a_0,a_1,a_2,b_0,b_1,b_2$. Put
\[
\Aff=\Span\{1,a_0,a_1,a_2,b_0,b_1,b_2\},\quad
P_0=a_0\wedge b_0,\quad Z_1=a_0\wedge b_1+a_1\wedge b_0,
\]
\[
z_0=a_0b_0,\quad z_1=a_0b_1+a_1b_0,\quad
z_2=a_0b_2+a_1b_1+a_2b_0+a_0a_1b_0b_1.
\]
Write $U=\langle a_0,a_1,b_0,b_1\rangle$, $\Omega=a_0\wedge a_1\wedge b_0\wedge b_1$, and $\delta_e=D_ez_2\in\langle a_0,b_0\rangle$ for an external direction $e$. The model is $g_i=L_iR_i$, $L_i,R_i\in V_i=\Aff+\Span\{g_0,\ldots,g_{i-1}\}$. Linear means that both factors are affine as functions. For a quadratic space $Q$, $V(Q)$ consists of the functions of degree at most two whose quadratic parts belong to $Q$.

Use as GIVEN [L4, Lemma 1] (normal form), [L4, Lemma 2] (quadratic spans), [L4, Theorem 12]; [L3, Lemma P (Lemma 6)], and [L3, Theorem 3]. The companion draft [L2, Lemma 2] proves the checked replacement argument, [L2, Lemma 3] proves the 32-element symmetry, and [L2, Lemma 5] proves the four-multiplier lemma. The human theorem of [L2] is now PROVED by ideal extraction and the one-high-signal obstruction; no open result from [L2] is used.

Here $g_0=a_0b_0$, $V_1=\Aff+\langle a_0b_0\rangle$, and $g_1,g_2,g_3,g_4$ are nonlinear.

\section{The theorem and the plan}
\begin{theorem}\label{thm:main}\status{PROVED}
No normalized five-AND circuit with exactly one linear gate computes $(z_0,z_1,z_2)$.
\end{theorem}
The quadratic branch reduces to [L2]. We then treat core and noncore cubic first gates, excluding both early top-bit placement and the final product. The full finite exclusion is independently verified below.

\begin{lemma}[Quadratic first gate]\label{lem:reduce}\status{PROVED}
If $\deg g_1\le2$, a five-gate circuit with one linear gate can be replaced by a five-gate circuit with exactly two linear gates that computes the same targets. The proved two-linear-gate theorem excludes this branch.
\end{lemma}
\begin{proof}
Write $g_1=fh$ with $f,h\in V_1$. By [L3, Lemma P (Lemma 6)], $(fh)_2=\pi+\rho$, where $\pi\in\langle P_0\rangle$ and $\rho=0$ or has rank two. Insert an affine-input gate with quadratic part $\rho$. If $\rho=0$ or $\rho=P_0$, choose any independent rank-two direction instead. The new two-linear-gate space contains $V_1+\langle g_1\rangle$ by [L4, Lemma 2]. Keep the old $g_2,g_3,g_4$ factors as functions, successively expressing them in the enlarged spaces. They remain nonlinear as functions, so the new circuit has exactly two linear gates. This is also the $m=1$ instance of [L2, Lemma 2].
\end{proof}
The companion [L2, Theorem 1] now supplies the required human exclusion, so the entire quadratic branch is PROVED.

\begin{lemma}[Cubic first-gate normal form]\label{lem:cubic}\status{PROVED}
Modulo $V_1$, every first nonlinear gate has a representative
\[
g_1\equiv(P_0+A)B\pmod{V_1},
\]
with $A,B$ homogeneous linear functions. It has degree three exactly when $B\notin\langle a_0,b_0\rangle$. Its cubic part is $P_0\wedge B$. The nonlinear factor planes modulo constants number 2,016; 1,920 give degree three and 96 give degree at most two.
\end{lemma}
\begin{proof}
The quadratic coefficients of the factors form a nonzero vector in the one-dimensional space $\langle P_0\rangle$. By exchanging factors or replacing one by their sum, assume one factor has quadratic part $P_0$ and the other is affine. Such plane operations change the product only by an available factor. Remove factor constants, also changing the product only by elements of $V_1$. We obtain $(P_0+A)B$.

The affine factor is a nonzero homogeneous linear function $B$. The cubic part is $P_0\wedge B$, which vanishes exactly for $B\in\langle a_0,b_0\rangle$. The factor plane has a unique nonzero affine vector modulo constants, namely $B$; its other vector is specified by $A$ modulo $\langle B\rangle$. There are 63 choices for $B$ and 32 choices for $A$ modulo $B$, giving $63\cdot32=2,016$ planes. Exactly 60 of the 63 choices lie outside $\langle a_0,b_0\rangle$, giving $60\cdot32=1,920$ cubic planes. Products cannot have degree four, since both quadratic parts lie in $\langle P_0\rangle$ and their wedge is zero.
\end{proof}

\begin{lemma}[The first nonlinear gate places no remaining output]\label{lem:first}\status{PROVED}
For every $g_1=fh$, $f,h\in V_1$, none of $z_1,z_2,z_1+z_2$ belongs to $V_2=V_1+\langle g_1\rangle$.
\end{lemma}
\begin{proof}
Lemma~\ref{lem:cubic} gives $\deg g_1\le3$, excluding $z_2$ and $z_1+z_2$, both of degree four. If $z_1\in V_2$, then $g_1\in z_1+V_1$ is quadratic. The product lemma gives $(g_1)_2=\pi+\rho$, with $\pi\in\langle P_0\rangle$ and $\rho$ zero or rank two. But $\rho$ would equal $Z_1$ or $Z_1+P_0$, both of rank four. This is impossible.
\end{proof}

\section{A cubic extension of the core coset argument}
Call a first-gate function core supported if all its terms of degree at least two use only $U=\langle a_0,a_1,b_0,b_1\rangle$. Free affine terms can then be removed so that the function depends only on $u=(a_0,a_1,b_0,b_1)$.

\begin{lemma}[Coset obstruction with a core cubic signal]\label{lem:corecoset}\status{PROVED}
Let $c$ be a core-supported cubic first gate from Lemma~\ref{lem:cubic}, and put $A=\Aff+\langle P_0,Z_1,c\rangle$. If $f,h\in A$, $g=fh$, and $W=A+\langle g\rangle$, then $z_2\notin W$ and no product from $W$ places $z_2$ in $W$.
\end{lemma}
\begin{proof}
Write $A=A_U\oplus\langle a_2,b_2\rangle$, removing the free external affine part of $c$. On $A_U$ define the bilinear form
\[
\beta(X,Y)=[a_0a_1b_0b_1](XY).
\]
Every function in $A_U$ has degree at most three, so $\beta(X,X)=0$ by $X^2=X$. Also $1,a_0,b_0$ belong to its radical. For $1$ this follows from the absence of a quartic term; for $a_0,b_0$ the only possible cubic term is $c_3=P_0\wedge B$, whose product with either cannot have a quartic term.

We first extend [L3, B3a (Lemma 15)]. Suppose $F,H\in A$, $v\in A+\langle fh\rangle$, and $FH+z_2+v$ has zero external derivatives and zero carry coefficient. Decompose
\[
f=\varphi+c_0,\quad h=\psi+d_0,\quad F=X+\ell,\quad H=Y+m,
\]
where the first terms belong to $A_U$ and the second terms are scalar external linear forms. Write $v=C+n+\sigma fh$. Comparing the two external linear coefficients gives
\[
\ell_iY+m_iX+\sigma(c_0)_i\psi+\sigma(d_0)_i\varphi
\equiv\lambda_i\pmod{\mathbb F_2},\qquad
\lambda_1=b_0,\quad\lambda_2=a_0.
\]
The external quadratic coefficient gives $\det(\ell,m)+\sigma\det(c_0,d_0)=0$, and the carry coefficient gives
\[
\beta(X,Y)+\sigma\beta(\varphi,\psi)=1.
\]
The coefficient matrix $M$ with columns $\ell,m,\sigma c_0,\sigma d_0$ has rank two, since the two $\lambda_i$ are independent modulo constants. Pass from $A_U$ to $A_U/\operatorname{rad}\beta$. The right sides vanish there because $a_0,b_0,1$ are radical. Thus the four resulting vectors lie in the two-dimensional coefficient kernel of $M$. Using its kernel basis, their pairings are its $2\times2$ minors times one common pairing. Complementary-minor duality [L3, Lemma 14] now gives exactly the cancellation of B3a: the coefficient of that common pairing is zero by the external quadratic equation. This contradicts the carry coefficient one. This replacement of quadratic parts by classes modulo $\operatorname{rad}\beta$ is the new step; it allows the cubic signal $c$.

The exact feedback conclusion of [L3, B3b (Lemma 16)] extends as well. Its derivative argument uses only constant external derivatives on the base space, membership of $D_e(fh)$ in that base, and the just-established B3a obstruction. These properties hold for $A$. More explicitly, if $k(g+F')=z_2+v$ with $k,F',v\in A$, differentiating gives the same identity as B3b. A nonzero external derivative of $k$ reduces it to a product of two base functions, already excluded. Otherwise $k$ depends only on $u$. The two external parts of $f,h$ must then be dependent; a direction killing both forces $k=\delta_e+\text{constant}$, and an independent direction would make a different nonconstant affine function vanish on its zero hyperplane. This is impossible, exactly as in the given B3b argument. The zero-external-parts case forces both independent functions $a_0,b_0$ to equal $k$ modulo constants and is also impossible.

Finally B3a excludes $z_2\in A+\langle g\rangle$. Any further product either has both factors in $A$, excluded by the extended B3a, or normalizes, including its output correction, to the feedback product excluded above. This proves the lemma.
\end{proof}

\begin{lemma}[Early $z_1$ forces a core cubic signal]\label{lem:earlycore}\status{PROVED}
Suppose $g_1$ is cubic and $z_1\in V_3$, after the next nonlinear gate $g_2$. Then $g_1$ is core supported. A feedback gate using $g_1$ cannot make this early placement at all.
\end{lemma}
\begin{proof}
Use $g_1\equiv(P_0+A)B\pmod{V_1}$ with $A,B$ homogeneous linear. First suppose $g_2$ has both factors in $V_1$. It must also be cubic; normalize it to $(P_0+C)D$. Cancellation of cubic terms gives $L=B+D\in\langle a_0,b_0\rangle$. With $M=A+C$, comparison of quadratic parts gives
\[
Z_1+\epsilon P_0=M\wedge B+C\wedge L.
\]
The left side has rank four and support $U$. Both summands must be nonzero rank-two forms, and [L4, Lemma 3] puts their supports in $U$. Hence $C,L,M,B\in U$, and therefore $A,B\in U$. The normalized $g_1$ is core supported; restoring elements of $V_1$ preserves that property.

If $g_2$ uses $g_1$, normalization gives $k(g_1+F)=z_1+v$ with $F,v\in V_1$. By [L2, Lemma 5], $k$ is one of the four complements of a fiber indicator (the constant 1 would make $g_1$ quadratic). On its three one-fibers, $g_1+F=z_1+v$. Project $A,B$ to the four free variables $(a_1,a_2,b_1,b_2)$, writing them $\bar A,\bar B$. The free quadratic part of this equality is $\bar A\wedge\bar B=0$. Cubicity implies $\bar B\ne0$, so $\bar A=\eta\bar B$.

Consequently $g_1$, on each fiber of $(a_0,b_0)$, is an affine function of those four free variables whose nonconstant coefficient vector changes only along the fixed one-dimensional direction $\bar B$. Adding $F$ does not change this rank bound. But the coefficient vector of $z_1+v$ changes in two independent directions, those of $a_1$ and $b_1$, over any three of the four fibers: any three points of $\mathbb F_2^2$ are affinely independent. The rank-one and rank-two changes cannot agree. This excludes the feedback case.
\end{proof}

\begin{proposition}[The core cubic branch]\label{prop:corebranch}\status{PROVED}
There is no five-gate target circuit whose first nonlinear gate is cubic and core supported.
\end{proposition}
\begin{proof}
Put $A=\Aff+\langle P_0,Z_1,g_1\rangle$. If $g_2$ places $z_1$, then $V_3=A$. The next product $g_3$ and the final gate contradict Lemma~\ref{lem:corecoset}. Otherwise $z_1\notin V_3$. Lemma~\ref{lem:corecoset}, applied to $g_2$, excludes $z_2$ and $z_1+z_2$ in $V_4$. The one-dimensional final extension therefore forces $z_1\in V_4$, so $V_4=A+\langle g_2\rangle$. Its last gate again contradicts the lemma.
\end{proof}

\section{Ordinary products in the noncore case}
In the following lemmas a first cubic gate can be replaced, modulo $V_1$, by $c=(P_0+L)B$ with $L$ affine and $B$ homogeneous linear. This changes none of the available spaces. Put $A_0=\Aff+\langle P_0,Z_1\rangle$ and $A=A_0+\langle c\rangle$.

\begin{lemma}[Two ordinary products cannot place $z_2$]\label{lem:ordinarypair}\status{PROVED}
For every cubic first gate $c$, there are no $F_1,F_2,F_3,F_4\in A$ with $F_1F_2+F_3F_4\in z_2+A$.
\end{lemma}
\begin{proof}
For core-supported $c$ this follows from Lemma~\ref{lem:corecoset}, applied to $g=F_1F_2$ and the further product $F_3F_4$. Suppose therefore that $c$ is noncore.

Write $F_i=a_i+t_i c$ with $a_i\in A_0$. Since $c^2=c$, the proposed identity would give $cH=z_2+v_0+a_1a_2+a_3a_4$ for some $H\in A_0$ and $v_0\in A_0$. Each product $a_i a_j$ has degree at most three: the quadratic parts lie in $\langle P_0,Z_1\rangle$, whose pairwise wedges vanish. Thus $cH$ has degree four with quartic part $\Omega$. But it is divisible as a Boolean function by $B$. The identity $B(cH)=cH$ forces $B\wedge\Omega=0$, so $B\in U$. Cubicity gives $B\notin\langle a_0,b_0\rangle$.

Noncore support now means that $L$ has a nonzero external part. Choose external linear coordinates $p,q$ so that $L=L_U+p$, where $L_U$ is affine on $U$. Then
\[
c=(P_0+L_U)B+pB.
\]
Write the external part of $a_i$ as $c_i p+d_iq$, and its core part as an affine function plus $w_iP_0+z_iZ_1$. The vectors $t,w,z,c_*,d$ are scalar four-vectors. On the hyperplane $B=j$ ($j=0,1$),
\[
F_i=x_i(j)+\alpha_i(j)p+d_iq,\qquad
\alpha(j)=c_*+jt,\quad
x_i(j)=(a_i)_U+jt_i(P_0+L_U).
\]
The quadratic coefficient vectors of the $x_i(j)$ are $w(j)=w+jt$ for $P_0$ and $z$ for $Z_1$.

Let $J$ interchange coordinates 1 and 2, and 3 and 4, and let the rows of the $2\times4$ matrix $M(j)$ be $(J\alpha(j))^T$ and $(Jd)^T$. Write the external part of $q_2$ as $sp+rq$, where $r,s$ are independent in $\langle a_0,b_0\rangle$. Coefficient comparison in the proposed identity gives
\[
M(j)x(j)\equiv(s,r)^T\pmod{\mathbb F_2},\qquad
M(j)JM(j)^T=0.
\]
The second equation is the zero $pq$ coefficient. The first implies $\rk M(j)=2$, because $r,s$ remain independent modulo constants on $B=j$.

The restrictions of $P_0,Z_1$ have independent quadratic parts on this hyperplane. Indeed $P_0$ cannot restrict to an affine function unless $B\in\supp P_0$, and $Z_1+\epsilon P_0$, of rank four, cannot lose its whole quadratic part on one hyperplane. Therefore
\[
M(j)w(j)=M(j)z=0.
\]
The rank-two isotropic row space of $M(j)$ has a two-dimensional isotropic kernel for the form $J$. Hence $w(j)^TJz=0$ for both values of $j$, and thus $w^TJz=t^TJz=0$.

For functions on the three-dimensional hyperplane $B=j$ define
\[
\beta_j(X,Y)=\sum_{u:B(u)=j}X(u)Y(u).
\]
On the restriction of $T=\Aff(U)+\langle P_0\rangle$, this is alternating, since every quadratic function on three bits has even parity. Its radical contains constants, $r$, and $s$: products of these with $P_0$ or affine functions have degree at most two on the hyperplane. Put $y_i(j)=x_i(j)+z_iZ_1$, so $y_i(j)\in T$ and $M(j)y(j)\equiv(s,r)^T$ modulo constants.

It follows that
\[
\beta_j(y_1,y_2)+\beta_j(y_3,y_4)=0.
\]
For completeness, subtract a particular solution with components in the radical span of $1,r,s$. The remaining vector lies in $\ker M(j)$. Complementary-minor duality [L3, Lemma 14] makes the displayed pairing the common kernel pairing times $\det_{12}M(j)+\det_{34}M(j)$, which is zero by isotropy. This is the same linear-algebra mechanism as B3a, now on each hyperplane.

Set $D=\sum_i(Jz)_i x_i(j)$. The relations $t^TJz=w^TJz=z^TJz=0$ show that $D$ is an affine function on $U$, independent of $j$. Expanding $x_i=y_i+z_iZ_1$ therefore gives
\[
\beta_j(x_1,x_2)+\beta_j(x_3,x_4)=\beta_j(Z_1,D).
\]
Here the pairing is also alternating on $T+\langle Z_1\rangle$, since its elements are quadratic on three bits. Summing over $j=0,1$ gives zero, because $Z_1D$ has degree at most three on four bits.

On the other hand, comparison at $p=q=0$ in the original identity makes that same sum the parity on all four core bits of $a_1b_1+\Omega$ plus a function of degree at most three. It is one. In general parity equals the coefficient of the full monomial, since every smaller monomial has an even number of ones. This contradiction proves the lemma.
\end{proof}

\section{Feedback multipliers modulo the low-bit ideal}
Let $I=(a_0b_0)$ as in [L2, Lemma 9]. Membership in $I$ means vanishing on the three fibers other than $(a_0,b_0)=(1,1)$.

\begin{lemma}[Noncore multiplier rigidity modulo $I$]\label{lem:noncoremult}\status{PROVED}
Suppose $c$ is a noncore cubic first gate. If $k\in A$, $R\in z_2+A$ and $kR=R$, then $k\equiv1\pmod I$.
\end{lemma}
\begin{proof}
Normalize $c=(P_0+L)B$ as above. Write $R=z_2+v+\sigma c$ and $k=u+\tau c$, with $u,v\in A_0$, $\sigma,\tau\in\mathbb F_2$. Put $w=1+k$, so $wR=0$.

First let $\tau=0$, so $k$ is quadratic. The degree-five part of $wR$ is exactly $\Omega\wedge u_{1,\mathrm{ext}}$. Indeed $u_2,v_2\in\langle P_0,Z_1\rangle$, $c_3=P_0\wedge B$, and $P_0\wedge Z_1=0$, so all other terms have degree at most four. Thus $u$ has no external linear part and $w$ depends only on $U$.

If $\sigma=0$, any external coefficient of $R$ gives $w(r+e)=0$ for some nonzero $r\in\langle a_0,b_0\rangle$ and constant $e$. If $\sigma=1$ and the external parts of $L,B$ are independent, $R$ has external quadratic coefficient one, forcing $w=0$. Otherwise those external parts span a single nonzero line (the zero-line case would make $c$ core). Choose external coordinates $p,q$ with $c$ independent of $q$. Again the $q$ coefficient gives $w(r+e)=0$.

This affine annihilation implies $w=(1+r+e)H$ with $H$ affine, since $w$ is quadratic. Hence $w_2$ has rank at most two. The only such nonzero form in $\langle P_0,Z_1\rangle$ is $P_0$; its factor $r$ then forces $H_1\in\langle a_0,b_0\rangle$. The zero-quadratic case does likewise. Therefore $w$ depends only on $a_0,b_0$.

On a fiber $(s,t)\ne(1,1)$, the quadratic part of $R$ in the four free variables is $a_1\wedge b_1+\sigma\bar L\wedge\bar B$, where bars project the homogeneous linear parts off $a_0,b_0$. It is nonzero. Otherwise $\sigma=1$ and the two projected factors span $\langle a_1,b_1\rangle$, forcing $L,B\in U$ and making $c$ core. Thus $wR=0$ forces $w=0$ on all three fibers, proving $w\in I$ in the quadratic case.

Now let $\tau=1$. The degree-six coefficient in $wR$ is the determinant of the external parts of $L,B$, times the nonzero monomial $a_2b_2\Omega$. All other contributions have degree at most five: $c_3$ already contains $P_0$, products with $P_0,Z_1$ have no degree five, and $c^2=c$. The two external parts must therefore be dependent. Since $c$ is noncore, they span a nonzero line. Choose external coordinates $p,q$ such that $c=\varphi(u)+pC(u)$ is independent of $q$, and write the external coefficients of $u$ as $u_p,u_q$.

If $u_q=1$, then $D_qk=1$. Differentiating $kR=R$ gives $R=kD_qR$, with $D_qR=r+e$ affine. This expresses $R$ as one product from $A$, contradicting Lemma~\ref{lem:ordinarypair}. Hence $u_q=0$. Comparing the $pq$ and $q$ coefficients of $wR=0$ now yields
\[
(C+u_p)(r+e)=0,\qquad (r+e)w=0.
\]
Here $r\in\langle a_0,b_0\rangle\setminus\{0\}$ and $e$ is constant.

If $B$ has no external part, then $L$ has external part $p$ and $C=B\in U\setminus\langle a_0,b_0\rangle$. The first displayed product of independent affine functions cannot vanish. Thus $B=B_U+p$, while $L=L_U+\eta p$. Using $B^2=B$, rewrite the same first gate exactly as
\[
c=(P_0+H)B,\qquad H=L_U+\eta B_U+\eta\in\Aff(U).
\]
Then $C=P_0+H$. The first displayed equation says that $P_0+H+u_p$ vanishes on $r=1+e$. On this hyperplane the restriction of $P_0$ is affine in the remaining low-bit coordinate; therefore $H$ has no $a_1,b_1$ coefficients. It depends only on $a_0,b_0$.

The function $C+u_p$ on those two low bits has quadratic coefficient one and vanishes on the two fibers $r=1+e$. Its parity is odd, so it is the indicator of a single fiber $j$ on $r=e$. Outside that hyperplane $w=0$. On the two fibers of the hyperplane $r=1+e$ we have $w=0$ and $C=u_p$; comparing their free-variable coefficient vectors in $w=1+u+c$ forces the $Z_1$ coefficient of $u$ to vanish. Indeed the $Z_1$ coefficient vector changes nontrivially between any two different low-bit fibers, while the affine coefficient vector and $CB_U$ have the same free part there.

Vanishing on either of those two fibers fixes the affine free part of $u$ to $u_p\bar B$. Hence on every fiber the free coefficient vector of $w$ is $(C+u_p)\bar B$, where $\bar B$ has external $p$ coefficient one. Thus $w$ is constant off the single fiber $j$, and is a nonconstant affine free-variable function with $p$ coefficient one on $j$.

On any fiber other than $(1,1)$, $c=CB$ is affine in the free variables. Consequently $R$ has free quadratic part exactly $a_1\wedge b_1$. If $j\ne(1,1)$, the product $wR$ has the nonzero cubic term $p\wedge a_1\wedge b_1$, impossible. Hence $j=(1,1)$. On the other three fibers $w$ is constant and $R$ has nonzero quadratic part, so $w=0$ there. Again $w\in I$, completing the proof.
\end{proof}

\section{The noncore branch and the theorem}
\begin{lemma}[Noncore intermediate-space obstruction]\label{lem:noncore}\status{PROVED}
Let $c$ be a noncore cubic first gate, $B=\Aff+\langle P_0,c\rangle$, $A=B+\langle Z_1\rangle$, and $g=fh$ a semantic nonlinear product with $f,h\in B$. Then $z_2\notin A+\langle g\rangle$, and no product from this enlarged space places $z_2$ in it.
\end{lemma}
\begin{proof}
Lemma~\ref{lem:ordinarypair} excludes $z_2\in A+\langle g\rangle$ and any further product with both factors in $A$, including a $g$ output correction: these are sums of at most two products from $A$.

Otherwise normalize the product and absorb the $g$ output correction to obtain $k(g+F)=z_2+v$ with $k,F,v\in A$. Lemma~\ref{lem:noncoremult} gives $k\equiv1\pmod I$. Write $c\equiv L\ell\pmod I$ using its first-gate normal form $(P_0+L)\ell$, with affine factors $L,\ell$. Put $\rho=(L\ell)_2$, which has rank at most two, and $Q=\langle P_0,\rho\rangle$.

Replacing $c$ by $L\ell$ in the factors of $g$ gives $f',h'\in V(Q)$ with $g\equiv f'h'\pmod I$. Replacing it likewise in $F,v$ gives $F',v'\in V(Q+\langle Z_1\rangle)$. The carry vanishes modulo $I$, so
\[
f'h'\equiv q_2+v'+F'\pmod I.
\]
Ideal extraction [L2, Lemma 9] forces a rank-at-most-two form in $q_2+Q+\langle Z_1\rangle$. Rank obstruction [L2, Lemma 10] gives minimum rank at least four (rank six if $\rho=0$). This contradiction proves the lemma.
\end{proof}

\begin{proof}[Proof of Theorem~\ref{thm:main}]
If $g_1$ is quadratic or affine, Lemma~\ref{lem:reduce} and [L2, Theorem 1] exclude the circuit. If it is cubic and core supported, Proposition~\ref{prop:corebranch} excludes it. Assume it is cubic and noncore.

By Lemma~\ref{lem:earlycore}, $z_1\notin V_3$. Put $A=\Aff+\langle P_0,g_1,Z_1\rangle$ and $g=g_2$. Lemma~\ref{lem:noncore} excludes both $z_2$ and $z_1+z_2$ from $V_4$: either placement, after making $z_1$ freely available in $A$, would place $z_2$ by adjoining at most one product to $A+\langle g\rangle$.

The one-dimensional final extension therefore forces $z_1\in V_4$, since otherwise the two nonzero target cosets would agree and put $z_1+z_2$ in $V_4$. As $z_1\notin V_3$, the third nonlinear gate places it and $V_4=A+\langle g\rangle$. The last gate now contradicts Lemma~\ref{lem:noncore} again. All cases are excluded.
\end{proof}

A quotient-only shortcut was tested and discarded: a relaxed four-gate circuit computes $(z_1,q_2)$ on the 48 points $a_0b_0=0$. Its explicit witness and independent simulation are retained in the checks. The proof above uses identities modulo $I$ only after the full-domain multiplier lemma has forced the multiplier to be one there; it does not rely on that false stronger relaxation.

\section{Computational checks and independent sanity check}
The standalone scripts use 64-bit ANF masks, exact binary elimination, and temporary truth tables for multiplication. The monomial-OR rule independently checks the multiplication transform. No SAT solver or third-party dependency is used. Every executable check runs a mutation or relaxed positive control.

The first-gate output is reduced modulo $V_1$ and its function space is reduced under the verified group of [L2, Lemma 3]. This yields 222 function-space orbits carrying all 2,016 first-plane multiplicities; 197 are cubic. These are function-space orbits, not the 249 factor-plane orbits in the supplied historical report. Redundant coordinates are removed as functions, including a first gate whose output already lies in $V_1$.

For each such space, enumerate every actual next nonlinear plane. Test placing $z_2$ or $z_1+z_2$ by $g_3$ using the exact screen described below. If $g_2$ places $z_1$, enumerate every $g_3$ plane and test the final gate. Otherwise enumerate all ways $g_3$ can place $z_1$, and test the final space $V_1+\langle g_1,g_2,z_1\rangle$. The early-target tests rule out placing $z_1$ only at $g_4$. All tests return no counterexample.

Here the screen works on any finite available space $B$, including spaces containing cubic signals. A new product either uses only factors in $B$, or normalizes to $k(g+F)$ with $k,F\in B$. For the latter, test all $k$ with a solution to $(1+k)(T+v)=0$, $v\in B$, and compare $T+kg$ with the complete linear image $kB$ modulo $B$. Changing the constant in $k$ accounts for the possible output term in $g$. Constant or dependent factors contribute only the available space, which is checked separately. Thus no gate is omitted. The last gate is relaxed to any product rather than only a nonlinear one.

\begin{center}\small
\begin{tabularx}{\textwidth}{@{}l X X@{}}\toprule
Item & Exact checked domain & Result\\\midrule
Lemma~\ref{lem:reduce} & All 2,016 first nonlinear planes & 96 quadratic-or-zero products; every replacement valid\\
Lemma~\ref{lem:cubic} & All 2,016 planes and their normalized factors & 1,920 cubic, 96 quadratic-or-zero; zero failures\\
Lemma~\ref{lem:first} & All first nonlinear planes, all three target cosets & 0 / 0 / 0 placements\\
Four-multiplier restriction & All $k\in V_1$, $R\in z_1+V_1$ & Included in the complete [L2, Lemma 5] check\\
Screen cross-check & All 1,740 distinct cubic first-output cosets; all next planes & 17,650,560 direct planes; exact agreement, including placing-plane counts\\
Lemma~\ref{lem:corecoset}, Proposition~\ref{prop:corebranch} & All 33 core cubic first spaces, all first and final products & Complete exact-space check; no failures\\
Lemma~\ref{lem:earlycore} & All 1,740 cubic output cosets and next planes & 48 placing pairs, all ordinary and core; no feedback placement\\
Lemma~\ref{lem:ordinarypair} & Every cubic first space; all pairs of ordinary products & 8,556,695 factor planes / 6,754,925 product cosets; no target pair\\
Lemma~\ref{lem:noncoremult} & All $k,R$ for 164 noncore cubic first spaces & 171,966,464 pairs; zero failures\\
Lemma~\ref{lem:noncore} & All 164 noncore cubic first spaces; no placement pre-screen & 1,426,210 next cosets; no final product\\
Full independent search & 222 function-space orbits, weight 2,016 & All early and final target branches empty\\
Positive controls & Ordinary and feedback relaxed targets & Detected in each executable check\\\bottomrule
\end{tabularx}\end{center}

With first-plane weights, the full search inspects 20,401,536 second-gate planes. It finds 72 $z_1$-placing second-gate planes, giving 3,080,448 candidate $(g_2,g_3)$ pairs in that branch. The later-placement branch has 1,542,672 pairs, for 4,623,120 in total. Every final test fails. The supplied report gives 4,644,000 using a different coordinate enumeration; that total is not reproduced here, and no field-by-field agreement is claimed. The exact spaces and all branches are nevertheless enumerated afresh; correctness of this run does not depend on historical counts. See the explicit coordinate-degeneracy discussion in \texttt{checks/README.md}.

\section{Status}
\begin{center}\begin{tabularx}{\textwidth}{@{}X l@{}}\toprule
Item & Status\\\midrule
Quadratic replacement implication (Lemma~\ref{lem:reduce}) & PROVED\\
Cubic normal form and first-placement exclusion (Lemmas~\ref{lem:cubic}, \ref{lem:first}) & PROVED\\
Four-multiplier restriction for the stated feedback branch & PROVED via [L2, Lemma 5]\\
Unconditional human exclusion of the quadratic branch & PROVED via [L2, Theorem 1]\\
Core cubic branch; early-$z_1$ core restriction & PROVED\\
Ordinary-product parity and noncore multiplier lemmas & PROVED\\
Noncore obstruction; Theorem~\ref{thm:main} & PROVED\\
Full independent normalized $\lambda=1$ enumeration & PASS (computation)\\\bottomrule
\end{tabularx}\end{center}
The computer result includes the quadratic and redundant branches directly and therefore does not depend on the human theorem in [L2]. The exhaustive search is an independent sanity check. The PROVED labels refer to the written arguments together with their passing exhaustive application-domain checks; human review is still pending.

\begin{thebibliography}{9}
\bibitem{L2} [L2] Companion draft, \emph{The two-linear-gate case of the low-half multiplier: a human proof and exhaustive checks}, \texttt{lambda2\_proof.tex}.
\bibitem{L3} [L3] Supplied \emph{A human proof of the three-linear-gate case for the $3\times3$ low-half multiplier}, \texttt{references/lambda3\_proof.tex}.
\bibitem{L4} [L4] Supplied \emph{A human proof of the four-linear-gate case for the $3\times3$ low-half multiplier}, \texttt{references/lambda4\_proof.tex}.
\end{thebibliography}
\end{document}
